\documentclass[11pt]{article}
\usepackage[margin=1in]{geometry}
\usepackage{enumitem}
% =============================================================================
% Preambles/header.tex — shared LaTeX/Beamer preamble for this project
%
% Use from a Beamer lecture by prepending:
%     \input{header}
% and compiling with TEXINPUTS=../Preambles:$TEXINPUTS (the /compile-latex
% skill does this automatically).
%
% PALETTE CONTRACT. Color names defined here MUST match the names in
% Quarto/theme-template.scss so Beamer and Quarto renderings use the same
% palette. The scripts/check-palette-sync.sh script enforces this.
%
% To customize: edit the HEX values below AND the matching $variable in
% Quarto/theme-template.scss, then run scripts/check-palette-sync.sh.
% =============================================================================

% -----------------------------------------------------------------------------
% Engine detection + encoding
%
% Our /compile-latex skill uses XeLaTeX, where inputenc is unnecessary and
% can produce encoding warnings. Only load inputenc under pdfTeX.
% -----------------------------------------------------------------------------
\usepackage{iftex}
\ifPDFTeX
  \usepackage[utf8]{inputenc}
\fi

\usepackage{amsmath, amssymb, amsthm}
\usepackage{booktabs}
\usepackage{graphicx}

% -----------------------------------------------------------------------------
% PALETTE — mirrors Quarto/theme-template.scss variable names.
% Load xcolor and define colors BEFORE anything that references them
% (notably hyperref, hypersetup, and the Beamer color assignments).
% -----------------------------------------------------------------------------
\usepackage{xcolor}

\definecolor{primary-blue}{HTML}{012169}      % headings, accents
\definecolor{primary-gold}{HTML}{B9975B}      % emphasis, borders
\definecolor{highlight-yellow}{HTML}{F2A900}  % markers, alerts
\definecolor{light-bg}{HTML}{E8EDF5}          % title / section backgrounds
\definecolor{jet}{HTML}{1A1A1A}               % body text

% Semantic colors (used in diagrams and annotations)
\definecolor{positive}{HTML}{15803D}          % good / observed / identified
\definecolor{negative}{HTML}{B91C1C}          % bad / problematic / bias
\definecolor{neutral}{HTML}{525252}           % reference / context

% Highlight accents (match .hi-* classes in the SCSS)
\definecolor{hi-slate}{HTML}{314F4F}
\definecolor{hi-green}{HTML}{2E8B57}
\definecolor{hi-red}{HTML}{C41E3A}

% Semantic aliases so skill templates and snippets can write
% \textcolor{accent}{...} without hard-coding "primary-blue".
\colorlet{accent}{primary-blue}
\colorlet{accent2}{primary-gold}

% -----------------------------------------------------------------------------
% Hyperref — loaded AFTER xcolor + palette so link colors resolve.
% -----------------------------------------------------------------------------
\usepackage{hyperref}
\hypersetup{colorlinks=true, linkcolor=primary-blue, urlcolor=primary-blue,
            citecolor=primary-blue, pdfborder={0 0 0}}

% -----------------------------------------------------------------------------
% Course code — set once per deck (after \input{header}, before \begin{document})
% via \coursecode{CS401}. Shown in the footer on every slide so a deck's course
% is identifiable without opening the title page. Empty by default.
% -----------------------------------------------------------------------------
\makeatletter
\newcommand{\coursecode}[1]{\gdef\@coursecodetext{#1}}
\gdef\@coursecodetext{}
\makeatother

% -----------------------------------------------------------------------------
% Beamer theme assignments (only apply under Beamer).
%
% We detect Beamer via \@ifclassloaded{beamer}, which is the documented,
% non-internal way. This must be wrapped in \makeatletter/\makeatother
% because \@ifclassloaded contains an @-catcode character.
% -----------------------------------------------------------------------------
\makeatletter
\@ifclassloaded{beamer}{%
  \usefonttheme{professionalfonts}
  \setbeamercolor{structure}{fg=primary-blue}
  \setbeamercolor{frametitle}{fg=primary-blue}
  \setbeamercolor{title}{fg=primary-blue}
  \setbeamercolor{subtitle}{fg=primary-gold}
  \setbeamercolor{author}{fg=primary-blue}
  \setbeamercolor{institute}{fg=neutral}
  \setbeamercolor{date}{fg=primary-gold}
  \setbeamercolor{itemize item}{fg=primary-blue}
  \setbeamercolor{itemize subitem}{fg=primary-gold}
  \setbeamercolor{alerted text}{fg=primary-gold}
  \setbeamercolor{block title}{fg=white, bg=primary-blue}
  \setbeamercolor{block body}{bg=light-bg}
  % Semantic block variants: block=definition/concept (blue), exampleblock=
  % worked example (green/positive), alertblock=key takeaway or caution (gold).
  % Keep to <=2 colored boxes per slide across ALL three variants (INV-7).
  \setbeamercolor{block title example}{fg=white, bg=positive}
  \setbeamercolor{block body example}{bg=light-bg}
  \setbeamercolor{block title alerted}{fg=white, bg=primary-gold}
  \setbeamercolor{block body alerted}{bg=light-bg}

  % Minimal footer: course code (left) + page X / N (right), no navigation clutter
  \setbeamertemplate{navigation symbols}{}
  \setbeamertemplate{footline}{%
    \color{neutral}\footnotesize\hspace{0.7em}\@coursecodetext\hfill
    \insertframenumber/\inserttotalframenumber\hspace{0.7em}\vspace{0.3em}}
}{}
\makeatother

% -----------------------------------------------------------------------------
% TikZ setup — libraries the snippet gallery uses
% -----------------------------------------------------------------------------
\usepackage{tikz}
\usetikzlibrary{arrows.meta, positioning, calc, decorations.pathreplacing,
                fit, shapes.geometric, backgrounds}

% Shared TikZ styles — snippets and hand-written diagrams can pick these up
% via `\begin{tikzpicture}[every node/.style={...}]` or by naming specific
% styles (e.g., `\node[dag-node] (X) at (0,0) {X};`).
\tikzset{
  % Node styles for causal DAGs and similar
  dag-node/.style = {
    draw=primary-blue, fill=light-bg, rounded corners=2pt,
    minimum width=1.2cm, minimum height=0.9cm, align=center, font=\small
  },
  decision-node/.style = {
    draw=primary-gold, fill=white, diamond, aspect=1.8,
    minimum width=1.8cm, minimum height=1.2cm, align=center, font=\small,
    inner sep=1pt
  },
  % Edge styles — observed vs counterfactual
  observed-edge/.style = {-{Latex[length=2.5mm]}, thick, draw=jet},
  counterfactual-edge/.style = {-{Latex[length=2.5mm]}, thick, draw=neutral, dashed},
  confound-edge/.style = {-{Latex[length=2.5mm]}, thick, draw=negative, dashed},
  % Data-point styles for plots
  observed-dot/.style = {fill=primary-blue, draw=primary-blue, circle, inner sep=1.2pt},
  counterfactual-dot/.style = {fill=white, draw=neutral, circle, inner sep=1.2pt}
}

% -----------------------------------------------------------------------------
% Convenience macros
% -----------------------------------------------------------------------------
\newcommand{\muted}[1]{\textcolor{neutral}{#1}}
\newcommand{\key}[1]{\textcolor{primary-gold}{\textbf{#1}}}
\newcommand{\good}[1]{\textcolor{positive}{#1}}
\newcommand{\bad}[1]{\textcolor{negative}{#1}}

% Standout frame macro: full-bleed dark background for section transitions.
% Beamer-only (uses \begin{frame}/\setbeamercolor/\usebeamerfont) -- do not
% call from an article-class file (e.g. Notes/<CODE>/*-notes.tex). \muted,
% \key, \good, \bad, and \coursecode above are plain xcolor/gdef and are
% safe to use from either class; verified when Notes/article-class support
% was added.
% Expand: \transitionslide{Your transition title}
\newcommand{\transitionslide}[1]{%
  \begingroup
  \setbeamercolor{background canvas}{bg=primary-blue}
  \begin{frame}[plain]
    \centering
    \vfill
    {\usebeamerfont{title}\color{white}\Large #1\par}
    \vfill
  \end{frame}
  \endgroup
}

% Section divider frame (Beamer-only), an alternative to \transitionslide with a
% darker two-line style: near-black (jet) background, a small white label line
% above a larger gold title, and a thin gold rule along the bottom edge.
% Expand: \sectiondivider{Section 2}{Pointers}
\newcommand{\sectiondivider}[2]{%
  \begingroup
  \setbeamercolor{background canvas}{bg=jet}
  \begin{frame}[plain]
    \vspace*{3.0cm}
    {\color{white}\ttfamily\large #1\par}
    \vspace{0.5cm}
    {\color{primary-gold}\LARGE #2\par}
    \begin{tikzpicture}[remember picture, overlay]
      \draw[primary-gold, line width=2.5pt]
        ($(current page.south west)+(0,0.1)$) -- ($(current page.south east)+(0,0.1)$);
    \end{tikzpicture}
  \end{frame}
  \endgroup
}

% -----------------------------------------------------------------------------
% Channel watermark — "Concepts That Click" (YouTube).
%
% Article-class files (CompetitiveExam Q/A sets, Notes, Assignments) get a
% light diagonal draft watermark on every page plus a centered footer naming
% the channel. Beamer decks get a subtle TikZ background watermark instead
% (the footline already carries the course code). Centralized here so every
% MCQ PDF is branded without per-file edits.
% -----------------------------------------------------------------------------
\makeatletter
\@ifclassloaded{article}{%
  \usepackage{draftwatermark}
  \SetWatermarkText{Concepts That Click}
  \SetWatermarkScale{1}
  \SetWatermarkFontSize{2cm}
  \SetWatermarkLightness{0.86}
  \SetWatermarkAngle{45}
  \usepackage{fancyhdr}
  \pagestyle{fancy}
  \fancyhf{}
  \fancyfoot[C]{\footnotesize\textcolor{neutral}{Concepts That Click~|~YouTube: \href{https://www.youtube.com/@concepts.thatclick}{@concepts.thatclick}}}
  \renewcommand{\headrulewidth}{0pt}
  \renewcommand{\footrulewidth}{0pt}
  \fancypagestyle{plain}{%
    \fancyhf{}%
    \fancyfoot[C]{\footnotesize\textcolor{neutral}{Concepts That Click~|~YouTube: \href{https://www.youtube.com/@concepts.thatclick}{@concepts.thatclick}}}%
    \renewcommand{\headrulewidth}{0pt}%
    \renewcommand{\footrulewidth}{0pt}%
  }
}{}
\@ifclassloaded{beamer}{%
  \setbeamertemplate{background}{%
    \begin{tikzpicture}[remember picture, overlay]
      \node[rotate=45, opacity=0.055, align=center]
        at (current page.center)
        {\Huge\textcolor{neutral}{Concepts That Click}};
    \end{tikzpicture}%
  }
}{}
\makeatother 
\coursecode{JKSSB}
\renewcommand{\thesection}{10.\arabic{section}}
\newtheorem{example}{Example}[section]
\newenvironment{definitionbox}[1]{%
  \par\noindent\textbf{Definition (#1).}\ \itshape
}{\par\vspace{0.3em}}
\title{Storage Devices --- Detailed Lecture Notes}
\author{JKSSB Computer Awareness}
\date{\today}

\begin{document}
\maketitle

\section{The storage hierarchy}
Storage is organized into levels according to how close it sits to the
processor and what job it does. The CPU works directly with primary storage;
secondary storage holds data permanently until it is needed and feeds it to
main memory; tertiary storage is an archival level used for large, rarely
accessed collections.

\begin{definitionbox}{Primary storage}
The directly CPU-accessible working memory of a computer, such as RAM, ROM,
and cache; it holds the data and instructions currently in use.
\end{definitionbox}

\begin{definitionbox}{Secondary storage}
Non-volatile storage that retains data when the power is off and is not
accessed directly by the CPU --- for example a hard disk, SSD, optical disc,
or USB flash drive.
\end{definitionbox}

\begin{definitionbox}{Tertiary storage}
An archival storage level for large volumes of infrequently used data,
typically robotic tape libraries whose media are loaded on demand.
\end{definitionbox}

\begin{center}
\begin{tabular}{p{0.2\textwidth}p{0.4\textwidth}p{0.3\textwidth}}
\toprule
\textbf{Level} & \textbf{Examples} & \textbf{Key property}\\
\midrule
Primary & RAM, ROM, cache & Directly accessible by the CPU\\
Secondary & HDD, SSD, optical disc, USB drive & Non-volatile; holds files\\
Tertiary & Robotic tape library & Archival; loaded on demand\\
\bottomrule
\end{tabular}
\end{center}

The CPU does not read a file straight from secondary storage. The operating
system brings the required data into main memory first, and the CPU then works
with that copy. This is why the hierarchy matters for an exam scenario: a
large program on disk is limited by the memory and storage that actually feed
it, not only by processor speed.

\section{Magnetic storage}
Magnetic storage records data by magnetizing a surface, and it is read back by
detecting those magnetic patterns. Hard disks and magnetic tape are the two
common magnetic media, but they differ sharply in how they are reached.

\subsection{Hard disk drive}
A hard disk drive (HDD) stores data on \emph{rotating platters} coated with a
magnetic material. A \emph{read/write head} mounted on an actuator arm floats
just above the platter surface and reads or writes as the platter spins. The
head moves inward and outward to reach different tracks. Because the platters
spin and the arm moves, an HDD is a \emph{mechanical} device with moving parts.

\begin{definitionbox}{Hard disk drive (HDD)}
Magnetic secondary storage that stores data on rotating platters read and
written by a moving read/write head; it is mechanical and contains moving
parts.
\end{definitionbox}

\subsection{Magnetic tape}
Magnetic tape stores data on a long, thin magnetic ribbon wound in a reel or
cartridge. The tape is read or written in order, position by position, so a
record far from the current position takes longer to reach. This makes tape a
\emph{sequential-access} medium. It is inexpensive and dense, which is why it
is used for backup and long-term archival rather than for routine access.

\begin{definitionbox}{Magnetic tape}
Magnetic storage on a long tape that is read and written in order; it is a
sequential-access medium used mainly for backup and archival.
\end{definitionbox}

\section{Optical storage}
Optical storage records data as \emph{pits and lands} --- tiny physical
markings --- on a disc and reads them with a \emph{laser}. The reflected light
pattern encodes the data. Optical discs are removable and non-volatile, so
they are convenient for distributing and archiving files.

The three common families are the compact disc (CD), the digital versatile
disc (DVD), and the Blu-ray disc. Blu-ray uses a shorter-wavelength laser and
tighter tracks, allowing it to store more than a DVD, which in turn stores
more than a CD.

\begin{definitionbox}{Optical storage}
Storage in which a laser reads data encoded as pits and lands on a disc, such
as a CD, DVD, or Blu-ray disc.
\end{definitionbox}

\begin{center}
\begin{tabular}{lll}
\toprule
\textbf{Medium} & \textbf{Approximate capacity (single layer)} &
\textbf{Reading method}\\
\midrule
CD & 700 MB & Laser\\
DVD & 4.7 GB & Laser (shorter wavelength)\\
Blu-ray & 25 GB & Laser (shortest wavelength)\\
\bottomrule
\end{tabular}
\end{center}

Order the capacities from smallest to largest as CD $<$ DVD $<$ Blu-ray.
A common test error is to swap the middle entries or to reverse the order
entirely; fix the sequence by remembering that CD is the oldest and smallest
and Blu-ray is the newest and largest.

\begin{example}
An item asks which optical disc holds the most data. Compare the standard
single-layer figures: CD $\approx$ 700 MB, DVD $\approx$ 4.7 GB, and Blu-ray
$\approx$ 25 GB. Since $25\ \text{GB} > 4.7\ \text{GB} > 700\ \text{MB}$, the
Blu-ray disc holds the most.
\end{example}

\section{Flash and solid-state storage}
Flash memory stores data in semiconductor cells that hold their state without
power. There are no rotating platters and no moving heads, so flash devices
are \emph{solid-state}. Common examples are the solid-state drive (SSD), the
USB flash drive, and the memory card.

\begin{definitionbox}{Flash memory}
Non-volatile semiconductor storage that keeps data in cells with no moving
parts; the technology behind SSDs, USB flash drives, and memory cards.
\end{definitionbox}

\begin{definitionbox}{Solid-state drive (SSD)}
A flash-based storage device with no moving parts. It is faster and more
shock-resistant than a mechanical hard disk and has a high transfer rate.
\end{definitionbox}

Because nothing spins and no head has to move, flash can start a transfer
almost immediately and is far less sensitive to being bumped or dropped.
Among a hard disk drive, an optical drive, magnetic tape, and an SSD, the
\key{SSD} has the highest transfer rate.

\section{Cloud storage}
Cloud storage keeps data on remote servers in a data centre and makes it
available over the Internet. The user sees a folder or an application, while
the provider manages the underlying hardware. Even though the storage feels
abstract, it ultimately rests on physical secondary storage --- arrays of
disks and solid-state drives --- on the provider's machines.

\begin{definitionbox}{Cloud storage}
Storage in which data is held on remote servers and accessed over the
Internet rather than on a locally attached device.
\end{definitionbox}

Cloud storage is not main memory and it is not the same as primary storage.
It is remote secondary storage reached through a network connection, useful
for backup, sharing, and access from several devices.

\section{Hard disk versus solid-state drive}
The comparison below is the most frequently tested near-neighbour pair in this
topic. The hard disk is magnetic and mechanical; the SSD is flash and
solid-state. Keep the technology and the consequences together.

\begin{center}
\begin{tabular}{p{0.2\textwidth}p{0.35\textwidth}p{0.35\textwidth}}
\toprule
\textbf{Aspect} & \textbf{Hard disk drive (HDD)} & \textbf{Solid-state drive (SSD)}\\
\midrule
Technology & Magnetic & Flash memory\\
Moving parts & Yes --- spinning platters, moving head & None\\
Speed & Usually slower transfer rate & Higher transfer rate (fastest here)\\
Shock resistance & More vulnerable to shock & More shock-resistant\\
Typical use & Bulk data at low cost per GB & Fast system and boot drive\\
\bottomrule
\end{tabular}
\end{center}

\begin{example}
A candidate reads ``the fastest storage device is the hard disk.'' That is
incorrect. The hard disk is magnetic and mechanical, with a spinning platter
and a moving head. The SSD is flash-based with no moving parts and has the
higher transfer rate. The correct choice is the SSD (KB PYQ-27).
\end{example}

\section{Access methods: random and sequential}
An access method describes how quickly a device can reach a stored location.

\begin{definitionbox}{Random (direct) access}
An access method in which any location can be reached in approximately equal
time; RAM and hard disks behave this way.
\end{definitionbox}

\begin{definitionbox}{Sequential access}
An access method in which locations are reached in order, so the time depends
on how far the target is from the current position; magnetic tape works this
way.
\end{definitionbox}

\begin{center}
\begin{tabular}{p{0.22\textwidth}p{0.3\textwidth}p{0.38\textwidth}}
\toprule
\textbf{Access method} & \textbf{Behavior} & \textbf{Examples}\\
\midrule
Random (direct) & Any location in about equal time & RAM, hard disk\\
Sequential & Reached in order & Magnetic tape\\
\bottomrule
\end{tabular}
\end{center}

Memory and disks are approximately random access because the electronics or
the head can be directed to a location without stepping through everything
before it. Tape is sequential because the ribbon must physically pass over the
read head in order.

\section{Non-volatility of storage}
Non-volatility is the property that stored contents survive when the power is
switched off. Storage devices are designed to be non-volatile so that files
remain available after shutdown. Primary storage such as RAM is
\emph{volatile}: its contents disappear when power is removed.

\begin{definitionbox}{Non-volatile storage}
Storage that retains its contents without power; hard disks, SSDs, optical
discs, and USB drives are non-volatile.
\end{definitionbox}

The one-line rule to carry into the exam is: RAM is volatile; storage is
non-volatile. This is exactly why unsaved work in RAM can be lost in a power
failure while a file already saved to disk remains.

\section{Exam retrieval and common confusions}
Use a two-direction recall drill. Given a device, state its technology and its
access method; given a technology or capacity, state the matching device.
\begin{itemize}
  \item HDD $\leftrightarrow$ magnetic, mechanical, random access.
  \item Magnetic tape $\leftrightarrow$ magnetic, sequential access, backup.
  \item CD / DVD / Blu-ray $\leftrightarrow$ optical, capacities 700 MB /
    4.7 GB / 25 GB (single layer).
  \item SSD / USB drive / memory card $\leftrightarrow$ flash, no moving
    parts; the SSD has the highest transfer rate.
  \item Cloud storage $\leftrightarrow$ remote servers over the Internet.
  \item RAM volatile versus storage non-volatile.
\end{itemize}

Three distinctions prevent most errors in this topic:
\begin{enumerate}
  \item Do not call the SSD magnetic or mechanical. The magnetic and
    mechanical device is the hard disk; the SSD is flash with no moving parts.
  \item Do not claim that the hard disk has the highest transfer rate; the
    SSD does (KB PYQ-27).
  \item Do not reverse the optical ordering. It is CD $<$ DVD $<$ Blu-ray
    for standard single-layer capacities, and do not confuse volatile RAM
    with non-volatile storage.
\end{enumerate}

\begin{example}
An item asks which device stores data without any moving parts. Eliminate the
hard disk (spinning platters and a moving head) and magnetic tape (a moving
ribbon). An optical drive also spins its disc. The flash-based SSD stores data
in semiconductor cells with no moving parts, so it is the answer.
\end{example}

\subsection*{Exercises}
\begin{enumerate}[label=\arabic*.]
  \item Distinguish primary, secondary, and tertiary storage and give one
    example of each.
  \item Explain why an SSD is described as having no moving parts, and why it
    is more shock-resistant than a hard disk.
  \item Order CD, DVD, and Blu-ray by single-layer capacity from smallest to
    largest, with the approximate figures.
  \item Classify RAM, a hard disk, and magnetic tape as random access or
    sequential access.
  \item A student writes: ``RAM is non-volatile, and the hard disk is the
    fastest storage device.'' Correct both errors.
\end{enumerate}

\subsection*{Solutions}
\begingroup\small
\begin{enumerate}[label=\arabic*.]
  \item Primary storage is directly accessible by the CPU (for example RAM).
    Secondary storage is non-volatile and holds data when power is off (for
    example a hard disk). Tertiary storage is archival storage loaded on
    demand (for example a robotic tape library).
  \item An SSD stores data in flash semiconductor cells, with no spinning
    platter and no moving read/write head. With no moving parts to be
    disturbed, it tolerates shock better than a hard disk and transfers data
    faster.
  \item CD $\approx$ 700 MB, then DVD $\approx$ 4.7 GB, then Blu-ray
    $\approx$ 25 GB, for single-layer media.
  \item RAM is random access; a hard disk is approximately random access;
    magnetic tape is sequential access.
  \item RAM is \emph{volatile}, not non-volatile: its contents are lost
    without power. The device with the highest transfer rate is the
    \emph{SSD}, not the hard disk; the hard disk is magnetic and mechanical
    with moving parts.
\end{enumerate}
\endgroup
\end{document}
